\chapter{Conclusion}

This thesis investigated the predictive value of sentiment extracted from financial filings for stock market forecasting. The analysis considered multiple dimensions, including alternative label definitions, different aggregation levels (company, sector, and industry), multiple model families, and temporal modeling using sequential deep learning architectures. Across these experiments, several consistent patterns emerged regarding when sentiment information is most useful and which modeling choices lead to reliable improvements.

A first important finding is that the usefulness of sentiment features depends strongly on the prediction task. Across nearly all experiments, sentiment information provided clearer benefits for classification problems than for regression. Improvements in directional prediction were observed for several model configurations, particularly when industry-adjusted labels or aggregated group modeling were employed. In contrast, regression models predicting numerical returns rarely benefited from sentiment features, and in many cases performance slightly degraded. This suggests that textual sentiment extracted from financial reports is more informative about the direction of price movements rather than their magnitude, which is likely influenced by additional market factors not captured by textual disclosures alone.

The choice of label definition also plays an important role. Industry-adjusted labels produced the most consistent improvements, particularly for linear models such as logistic regression. Removing broader industry movements appears to isolate firm-specific signals contained in financial disclosures, allowing sentiment features to contribute more effectively. Sector-adjusted labels showed weaker and less consistent effects, while raw labels exhibited moderate improvements depending on the model. These results indicate that sentiment information is most useful when broader systematic movements are controlled for, highlighting the importance of careful label construction in financial prediction tasks.

The aggregation level of the modeling framework further influences the effectiveness of sentiment features. Industry-level experiments yielded the most consistent classification improvements, followed by sector-level modeling, while company-level results were more heterogeneous. This pattern suggests that sentiment signals extracted from financial reports capture shared dynamics within industries, such as common regulatory environments, demand cycles, or macroeconomic exposure. At the company level, the smaller number of observations introduces additional variance, reducing statistical reliability despite occasionally large individual improvements. These findings highlight a trade-off between model specialization and statistical robustness, with industry-level aggregation providing the best balance.

Model choice also significantly affects performance. Linear models benefited most consistently from the inclusion of sentiment features, particularly in classification settings. These models have limited capacity to capture nonlinear relationships from structured financial variables, making them more receptive to additional textual information. In contrast, tree-based ensemble models such as Random Forest and XGBoost often exhibited smaller improvements or slight degradations. This behaviour suggests that these models already capture complex interactions within structured inputs, leaving less room for incremental gains from sentiment features.

The time-series experiments further demonstrated that temporal modeling enhances the usefulness of sentiment information, particularly for classification tasks. The LSTM architecture provided the most stable improvements, especially at the industry level and for intermediate sequence lengths at the sector level. At the company level, shorter sequences produced the strongest results, indicating that recent filing sentiment contains the most informative signal and that the predictive value of sentiment decays over time. The PatchTST transformer architecture produced mixed results but showed advantages for regression tasks with short sequences, suggesting that attention-based models may better capture short-term dependencies relevant for return prediction.

Taken together, the results indicate that sentiment extracted from financial filings contains meaningful predictive information, particularly for directional forecasting. The most consistent gains are obtained when using industry-adjusted labels, industry-level aggregation, and temporal modeling with moderate sequence lengths. Improvements are strongest for classification tasks, while regression performance remains limited and model-dependent. These findings support the conclusion that textual sentiment is more effective for identifying market direction rather than estimating precise return magnitudes.

\section{Future Work}

Several directions for future research emerge from the findings of this study. A first promising extension involves improving the representation of textual information beyond aggregated sentiment scores. While this work relies on document-level sentiment extracted from financial reports, richer representations such as section-level sentiment, topic distributions, or discourse-aware features may capture more nuanced signals. For example, separating management discussion, risk disclosures, and forward-looking statements could allow models to differentiate between short-term operational changes and long-term strategic outlooks.

Another direction concerns the integration of additional textual sources. The present study focuses on formal financial filings, which are typically released at relatively low frequency. Incorporating higher-frequency sources such as earnings call transcripts, press releases, or financial news could provide more timely sentiment signals and improve predictive performance, particularly for short-horizon forecasting. Combining multiple textual modalities may also help reduce noise and increase robustness across companies and sectors.

Future work may also explore more advanced temporal modeling strategies. Although sequence-based models such as LSTM and PatchTST capture dependencies across filings, they treat each time step uniformly. Attention mechanisms over time, adaptive weighting of recent filings, or decay-based temporal aggregation could better model the diminishing relevance of older disclosures. Additionally, hierarchical temporal architectures that jointly model company-, sector-, and industry-level dynamics may further improve predictive performance.

Another potential extension involves improving the treatment of cross-sectional dependencies. The current experiments evaluate grouping strategies independently at the company, sector, and industry levels. Graph-based approaches, such as graph neural networks constructed from supply chain relationships, industry similarity, or return correlations, could allow sentiment signals to propagate between related entities. This may help stabilize company-level predictions and improve performance for firms with limited historical data.

Finally, future research could investigate alternative evaluation settings and prediction targets. The present work focuses on short-term return direction and magnitude; however, sentiment information may be more informative for volatility prediction, abnormal return detection, or event-driven strategies. Additionally, incorporating transaction costs and portfolio construction into the evaluation framework would provide a more realistic assessment of the economic value of sentiment-based predictions.

Overall, these extensions could further clarify the role of textual sentiment in financial forecasting and help bridge the gap between statistically significant predictive improvements and practical investment applications.